% problem_id: S001-lemma-cohomology-projective-space-over-ring
% source_id: S001 (Stacks Project)
% source_file: coherent.tex
% statement_locator: coherent.tex lines 1602--1622
% proof_locator: coherent.tex lines 1624--1840
% env: lemma
% id_note: label 在本来源内唯一
% pairing: tight (verified)
% status: complete (statement + author proof verbatim + reason + locators)
%
\begin{lemma}
\label{lemma-cohomology-projective-space-over-ring}
\begin{reference}
\cite[III Proposition 2.1.12]{EGA}
\end{reference}
Let $R$ be a ring.
Let $n \geq 0$ be an integer.
We have
$$
H^q(\mathbf{P}^n, \mathcal{O}_{\mathbf{P}^n_R}(d)) =
\left\{
\begin{matrix}
(R[T_0, \ldots, T_n])_d & \text{if} & q = 0,\ d \geq 0 \\
\left(\frac{1}{T_0 \ldots T_n} R[\frac{1}{T_0}, \ldots, \frac{1}{T_n}]\right)_d
& \text{if} & q = n,\ d < 0 \\
0 & \text{if} & q, n, d\text{ not as above}
\end{matrix}
\right.
$$
as $R$-modules.
\end{lemma}

\begin{proof}
We will use the standard affine open covering
$$
\mathcal{U} : \mathbf{P}^n_R = \bigcup\nolimits_{i = 0}^n D_{+}(T_i)
$$
to compute the cohomology using the {\v C}ech complex.
This is permissible by Lemma \ref{lemma-cech-cohomology-quasi-coherent}
since any intersection of finitely many affine $D_{+}(T_i)$ is also a
standard affine open (see
Constructions, Section \ref{constructions-section-proj}).
In fact, we can use the alternating or ordered {\v C}ech complex according to
Cohomology, Lemmas \ref{cohomology-lemma-ordered-alternating} and
\ref{cohomology-lemma-alternating-usual}.

\medskip\noindent
The ordering we will use on $\{0, \ldots, n\}$ is the usual one.
Hence the complex we are looking at has terms
$$
\check{\mathcal{C}}_{ord}^p(\mathcal{U}, \mathcal{O}_{\mathbf{P}_R}(d))
=
\bigoplus\nolimits_{i_0 < \ldots < i_p}
(R[T_0, \ldots, T_n, \frac{1}{T_{i_0} \ldots T_{i_p}}])_d
$$
Moreover, the maps are given by the usual formula
$$
d(s)_{i_0 \ldots i_{p + 1}} =
\sum\nolimits_{j = 0}^{p + 1} (-1)^j s_{i_0 \ldots \hat i_j \ldots i_{p + 1}}
$$
see Cohomology, Section \ref{cohomology-section-alternating-cech}.
Note that each term of this complex has a natural
$\mathbf{Z}^{n + 1}$-grading. Namely, we get this by declaring a monomial
$T_0^{e_0} \ldots T_n^{e_n}$ to be homogeneous with weight
$(e_0, \ldots, e_n) \in \mathbf{Z}^{n + 1}$. It is clear that the differential
given above respects the grading. In a formula we have
$$
\check{\mathcal{C}}_{ord}^\bullet(\mathcal{U}, \mathcal{O}_{\mathbf{P}_R}(d))
=
\bigoplus\nolimits_{\vec{e} \in \mathbf{Z}^{n + 1}}
\check{\mathcal{C}}^\bullet(\vec{e})
$$
where not all summands on the right hand side occur (see below).
Hence in order to compute the cohomology
modules of the complex it suffices to compute the cohomology of the graded
pieces and take the direct sum at the end.

\medskip\noindent
Fix $\vec{e} = (e_0, \ldots, e_n) \in \mathbf{Z}^{n + 1}$. In order for this
weight to occur in the complex above we need to assume
$e_0 + \ldots + e_n = d$ (if not then it occurs for a different twist of
the structure sheaf of course). Assuming this, set
$$
NEG(\vec{e}) = \{i \in \{0, \ldots, n\} \mid e_i < 0\}.
$$
With this notation the weight $\vec{e}$ summand
$\check{\mathcal{C}}^\bullet(\vec{e})$ of the {\v C}ech complex above has
the following terms
$$
\check{\mathcal{C}}^p(\vec{e})
=
\bigoplus\nolimits_{i_0 < \ldots < i_p,
\ NEG(\vec{e}) \subset \{i_0, \ldots, i_p\}}
R \cdot T_0^{e_0} \ldots T_n^{e_n}
$$
In other words, the terms corresponding to $i_0 < \ldots < i_p$ such
that $NEG(\vec{e})$ is not contained in $\{i_0 \ldots i_p\}$ are zero.
The differential of the complex $\check{\mathcal{C}}^\bullet(\vec{e})$
is still given by the exact same formula as above.

\medskip\noindent
Suppose that $NEG(\vec{e}) = \{0, \ldots, n\}$, i.e., that all
exponents $e_i$ are negative.
In this case the complex $\check{\mathcal{C}}^\bullet(\vec{e})$ has
only one term, namely $\check{\mathcal{C}}^n(\vec{e}) =
R \cdot \frac{1}{T_0^{-e_0} \ldots T_n^{-e_n}}$. Hence in this
case
$$
H^q(\check{\mathcal{C}}^\bullet(\vec{e})) =
\left\{
\begin{matrix}
R \cdot \frac{1}{T_0^{-e_0} \ldots T_n^{-e_n}} & \text{if} & q = n \\
0 & \text{if} & \text{else}
\end{matrix}
\right.
$$
The direct sum of all of these terms clearly gives the value
$$
\left(\frac{1}{T_0 \ldots T_n} R[\frac{1}{T_0}, \ldots, \frac{1}{T_n}]\right)_d
$$
in degree $n$ as stated in the lemma. Moreover these terms do not contribute
to cohomology in other degrees (also in accordance with the statement of the
lemma).

\medskip\noindent
Assume $NEG(\vec{e}) = \emptyset$. In this case the complex
$\check{\mathcal{C}}^\bullet(\vec{e})$ has a summand $R$ corresponding
to all $i_0 < \ldots < i_p$.
Let us compare the complex $\check{\mathcal{C}}^\bullet(\vec{e})$
to another complex. Namely, consider the affine open covering
$$
\mathcal{V} : \Spec(R) = \bigcup\nolimits_{i \in \{0, \ldots, n\}} V_i
$$
where $V_i = \Spec(R)$ for all $i$. Consider the alternating
{\v C}ech complex
$$
\check{\mathcal{C}}_{ord}^\bullet(\mathcal{V}, \mathcal{O}_{\Spec(R)})
$$
By the same reasoning as above this computes the cohomology of the
structure sheaf on $\Spec(R)$. Hence we see that
$H^p(
\check{\mathcal{C}}_{ord}^\bullet(\mathcal{V}, \mathcal{O}_{\Spec(R)})
) = R$ if $p = 0$ and is $0$ whenever $p > 0$.
For these facts, see
Lemma \ref{lemma-cech-cohomology-quasi-coherent-trivial} and its proof.
Note that also
$\check{\mathcal{C}}_{ord}^\bullet(\mathcal{V}, \mathcal{O}_{\Spec(R)})$
has a summand $R$ for every $i_0 < \ldots < i_p$ and has exactly the same
differential as $\check{\mathcal{C}}^\bullet(\vec{e})$. In other words
these complexes are isomorphic complexes and hence have the same cohomology.
We conclude that
$$
H^q(\check{\mathcal{C}}^\bullet(\vec{e})) =
\left\{
\begin{matrix}
R \cdot T_0^{e_0} \ldots T_n^{e_n} & \text{if} & q = 0 \\
0 & \text{if} & \text{else}
\end{matrix}
\right.
$$
in the case that $NEG(\vec{e}) = \emptyset$.
The direct sum of all of these terms clearly gives the value
$$
(R[T_0, \ldots, T_n])_d
$$
in degree $0$ as stated in the lemma. Moreover these terms do not contribute
to cohomology in other degrees (also in accordance with the statement of the
lemma).

\medskip\noindent
To finish the proof of the lemma we have to show that the complexes
$\check{\mathcal{C}}^\bullet(\vec{e})$ are acyclic when
$NEG(\vec{e})$ is neither empty nor equal to $\{0, \ldots, n\}$.
Pick an index $i_{\text{fix}} \not \in NEG(\vec{e})$ (such an index exists).
Consider the map
$$
h :
\check{\mathcal{C}}^{p + 1}(\vec{e})
\to
\check{\mathcal{C}}^p(\vec{e})
$$
given by the rule that for $i_0 < \ldots < i_p$ we have
$$
h(s)_{i_0 \ldots i_p} =
\left\{
\begin{matrix}
0 & \text{if} & p \not \in \{0, \ldots, n - 1\} \\
0 & \text{if} & i_{\text{fix}} \in \{i_0, \ldots, i_p\} \\
s_{i_{\text{fix}} i_0 \ldots i_p} & \text{if} & i_{\text{fix}} < i_0 \\
(-1)^a s_{i_0 \ldots i_{a - 1} i_{\text{fix}} i_a \ldots i_p} &
\text{if} & i_{a - 1} < i_{\text{fix}} < i_a \\
(-1)^{p + 1} s_{i_0 \ldots i_p i_{\text{fix}}} &
\text{if} & i_p < i_{\text{fix}}
\end{matrix}
\right.
$$
Please compare with the proof of
Lemma \ref{lemma-cech-cohomology-quasi-coherent-trivial}.
This makes sense because we have
$$
NEG(\vec{e}) \subset \{i_0, \ldots, i_p\}
\Leftrightarrow
NEG(\vec{e}) \subset \{i_{\text{fix}}, i_0, \ldots, i_p\}
$$
The exact same (combinatorial) computation\footnote{
For example, suppose that $i_0 < \ldots < i_p$ is such that
$i_{\text{fix}} \not \in \{i_0, \ldots, i_p\}$ and that
$i_{a - 1} < i_{\text{fix}} < i_a$ for some $1 \leq a \leq p$. Then we have
\begin{align*}
&
(dh + hd)(s)_{i_0 \ldots i_p} \\
& =
\sum\nolimits_{j = 0}^p
(-1)^j
h(s)_{i_0 \ldots \hat i_j \ldots i_p}
+
(-1)^a d(s)_{i_0 \ldots i_{a - 1} i_{\text{fix}} i_a \ldots i_p}\\
& =
\sum\nolimits_{j = 0}^{a - 1}
(-1)^{j + a - 1}
s_{i_0 \ldots \hat i_j \ldots i_{a - 1} i_{\text{fix}} i_a \ldots i_p} +
\sum\nolimits_{j = a}^p
(-1)^{j + a}
s_{i_0 \ldots i_{a - 1} i_{\text{fix}} i_a \ldots \hat i_j \ldots i_p}
+ \\
&
\sum\nolimits_{j = 0}^{a - 1}
(-1)^{a + j}
s_{i_0 \ldots \hat i_j \ldots i_{a - 1} i_{\text{fix}} i_a \ldots i_p} +
(-1)^{2a} s_{i_0 \ldots i_p} +
\sum\nolimits_{j = a}^p
(-1)^{a + j + 1}
s_{i_0 \ldots i_{a - 1} i_{\text{fix}} i_a \ldots \hat i_j \ldots i_p}
\\
& =
s_{i_0 \ldots i_p}
\end{align*}
as desired. The other cases are similar.}
as in the proof of Lemma \ref{lemma-cech-cohomology-quasi-coherent-trivial}
shows that
$$
(hd + dh)(s)_{i_0 \ldots i_p}
=
s_{i_0 \ldots i_p}
$$
Hence we see that the identity map of the complex
$\check{\mathcal{C}}^\bullet(\vec{e})$ is homotopic to zero
which implies that it is acyclic.
\end{proof}
